% ============================================================================
% 06-discussao.tex — Discussão: raio-x da arquitetura (módulo 7/8)
% ============================================================================
\chapter{Discussão}

\section{Raio-X completo da arquitetura que conduziu ao resultado}

A Figura \ref{fig:raiox} apresenta as seis camadas da arquitetura
responsáveis pelo resultado: cada camada contribui com função específica e
não opcional no ciclo Perceber$\rightarrow$Especificar$\rightarrow$Delegar
$\rightarrow$Executar$\rightarrow$Verificar$\rightarrow$Refletir.

% ============================================================================
% fig-raiox.tex — Raio-X da arquitetura (Figura 3): camadas C0–C5
% ============================================================================
\begin{figure}[htbp]
\centering
\begin{tikzpicture}[
  node distance=0.55cm,
  camada/.style={rectangle, rounded corners=3pt, minimum width=13.5cm,
                 minimum height=1.15cm, align=left, font=\footnotesize,
                 draw=black, thick}
]
  \node[camada, fill=cororo!12, text width=13cm]
    (c0) {\textbf{C0 -- Infraestrutura}: host opencode (sessão 1h15m) + 6 MCP servers + CLI \texttt{opencode run} + pdflatex/abntex2};
  \node[camada, fill=corverde!10, text width=13cm, below=of c0]
    (c1) {\textbf{C1 -- Catálogo e roteamento}: ModelRouter (41 catálogo + 5 free curados R500); providers go/zen/litert/runai/opencode; \texttt{route\_free(task\_type)}};
  \node[camada, fill=coramarelo!10, text width=13cm, below=of c1]
    (c2) {\textbf{C2 -- Modelos free reais}: mimo (0,987) $>$ big-pickle (0,923) $>$ nemotron (0,709) $>$ muse-1.3 (0,324) $>$ muse-1.2 (0,312); locais descartados por latência};
  \node[camada, fill=corverde!10, text width=13cm, below=of c2]
    (c3) {\textbf{C3 -- Scaffold MASWOS (16 estágios)}: condição C2 decisiva 0\%$\rightarrow$100\% (1º trial); funcionamento análogo a CoT/Reflexion};
  \node[camada, fill=cororo!12, text width=13cm, below=of c3]
    (c4) {\textbf{C4 -- Verificação formal}: RealIMOSolver anti-\textit{leak} (2/4, média 4,5/7); GradingHead; TestRealPipeContract (deep=True)};
  \node[camada, fill=corvermelho!8, text width=13cm, below=of c4]
    (c5) {\textbf{C5 -- Governança e memória}: MetaBus; Trust Engine; EvolutionRegistry (361 ciclos); honest-critic; anti-overclaim (CORRIGENDUM)};

  \node[font=\scriptsize, color=corcinza, right=1.1cm of c0] (s0) {R498};
  \node[font=\scriptsize, color=corcinza, right=1.1cm of c1] (s1) {R500};
  \node[font=\scriptsize, color=corcinza, right=1.1cm of c2] (s2) {R499+réplicas};
  \node[font=\scriptsize, color=corcinza, right=1.1cm of c3] (s3) {R499};
  \node[font=\scriptsize, color=corcinza, right=1.1cm of c4] (s4) {R499/R498};
  \node[font=\scriptsize, color=corcinza, right=1.1cm of c5] (s5) {todo o ciclo};
\end{tikzpicture}
\caption{Raio-X da arquitetura do OpenCode Ecosystem Core que conduziu ao
resultado: camadas C0--C5 com os ciclos evolutivos correspondentes (R497--
R500). O fluxo de dados atravessa todas as camadas; a governança (C5) audita
as demais.}
\label{fig:raiox}
\end{figure}

\subsection*{Camada C0 --- Infraestrutura (\textit{host} e MCP)}
Seis servidores MCP persistentes (metacognitive-interconnect, scanners,
pypi-search, colibri, litert-lm, antigravity-bridge) + \textit{host} CLI
opencode (1h15m de sessão no momento da replicação) + \texttt{pdflatex} e
abntex2. A verificação por \texttt{ps} mostrou que os MCP não importam
\texttt{model\_router}: a atualização R500 não exige reinício deles; o
\textit{host} CLI, porém, acumula estado e é a hipótese dominante para o
\texttt{big-pickle} vazio na replicação.

\subsection*{Camada C1 --- Catálogo e roteamento}
\texttt{ModelRouter} com 41 modelos catalogados + 5 free curados (R500);
providers \texttt{opencode-go}, \texttt{opencode-zen}, \texttt{litert-lm},
\texttt{runai}, \texttt{opencode}. Sem \texttt{prefer\_free}, perfis clássicos
permanecem; com \texttt{prefer\_free}$\wedge$tarefa acadêmica/math/reasoning/
writing, a curadoria R499 decide (mimo $>$ big-pickle $>$ nemotron $>$
muse).

\subsection*{Camada C2 --- Modelos gratuitos reais}
mimo, big-pickle, nemotron, muse 1.2/1.3 (todos via CLI \texttt{opencode
run}); locais descartados por latência (llama3.2: 30--50~s/estágio $\times$ 16
estágios $>$ 10~min). O rótulo ``free'' não implica disponibilidade real:
deepseek (saldo), gemini (erro), gpt-4o-mini/claude-haiku (vazio).

\subsection*{Camada C3 --- \textit{Scaffold} de raciocínio}
O template MASWOS de 16 estágios (condição C2) foi o fator decisivo: elevou a
acurácia de 0\% para 100\% (primeiro \textit{trial}, modelos que resolveram).
Funcionalmente análogo a \textit{Chain-of-Thought} \cite{wei_cot}, porém com
etapas obrigatórias de verificação e conformidade de saída, aproximando o
padrão \textit{Reflexion} \cite{shinn_reflexion} quanto à autoavaliação.

\subsection*{Camada C4 --- Verificação formal}
O \texttt{RealIMOSolver} eliminou o \textit{leak} do \textit{harness} original
(R442) e classificou 2/4 problemas corretamente (média 4,5/7), expondo
ambiguidades de paráfrase (\texttt{N=1} vs \texttt{N=3}; máximo local com
extremos) e sensibilidade de serialização ($2^{n-1}$ vs $2^{(n-1)}$) — lições
para \textit{grading heads} automáticos. O \texttt{TestRealPipeContract}
(R498) garante que o pipeline de diagnóstico roda \texttt{deep=True} em
documentos reais.

\subsection*{Camada C5 --- Memória metacognitiva e governança}
MetaBus (lições compartilhadas), Trust Engine, Registro de Evolução (361
ciclos), banca simulada honesta (sem nota), política anti-
\textit{overclaim} (CORRIGENDUM). A banca simulada MASWOS local não concluiu
(abortada por latência); a decisão de não declarar nota é registrada como
prática.

\section{Agentes e processos multiagentes de melhor desempenho}

Entre os 206 agentes do catálogo, os de maior contribuição para este estudo
foram: (i) \textbf{honest-critic-agent} (R142) — bloqueou a alegação
``pronto como/melhor que superhuman'' (\textit{publishable=False}),
convergindo com a evidência; (ii) \textbf{agentes acadêmicos MASWOS} (33
agentes, A01--A16) — o \textit{scaffold} delegado via \texttt{delegate\_fn}
foi decisivo para o acerto dos LLMs na C2; (iii) \textbf{coder/test-engineer}
— implementação e testes R500 (46 verdes, 1,97~s); (iv) \textbf{agent de
curadoria landscape} — inserção de \texttt{deepmind/acme} (R497); (v)
\textbf{RealIMOSolver} (agente de verificação determinística) — sem ele, o
\textit{harness} ecoaria a resposta (\textit{leak}); (vi) \textbf{banca
rigorosa MASWOS + honest-critic} — gate editorial anti-
\textit{overclaim}. O processo A2A (Blackboard + MetaBus) assegurou que as
lições do R499 (leak, fake reasoning, estocasticidade) fossem consumidas no
R500.

\section{Comparação com a literatura}

O efeito de \textit{scaffold} observado (0\%$\rightarrow$100\%) é consistente
com \textit{Chain-of-Thought} \cite{wei_cot}, \textit{Reflexion}
\cite{shinn_reflexion} e \textit{ReAct} \cite{yao_react}: a estrutura
explícita de passos e autoavaliação importa mais que o modelo, para problemas
curtos de olimpíada. A indisponibilidade real de ``free'' (saldo/erro/vazio)
reforça o alerta de \cite{schaeffer_mirage}: métricas e condições de
\textit{deployment} mudam a conclusão. O \textit{score} composto (acurácia,
velocidade, conformidade, disponibilidade) é uma alternativa operacional às
metrificações binárias de \textit{benchmark}; a replicação (n=3) corrige o
viés de \textit{um trial} observado em relatórios como \cite{openai_gpt4}.

\section{Limitações}

(i) $n$ pequeno (4 problemas de \textit{harness}; 1 problema no \textit{run}
de benchmark por condição; n=3 máximo na replicação); (ii) instabilidade de
latência (timeout 110~s no mimo); (iii) \textit{host} sem reinício antes da
replicação — \texttt{big-pickle} não confirmado pós-R500; (iv) validação
externa ausente (nenhum periódico/editora revisou o manuscrito); (v) abntex2
instalado em \texttt{TEXMFHOME} (não no TeX Live do sistema); (vi) DOIs
validados por HTTP 200 em \texttt{doi.org}, não por revisão editorial
cruzada. Nenhuma dessas limitações invalida a evidência reportada; todas
delas impedem rótulos de excelência.